\chapter{Metodología y diseño experimental}
\label{ch:metodologia}

Este capítulo describe la metodología completa del estudio: los datos
utilizados, el esquema de backtesting walk-forward, las carteras evaluadas
y las métricas de rendimiento. El objetivo es que el experimento sea
completamente reproducible con los datos y el código proporcionados.

%--------------------------------------------------------------------
\section{Datos}
%--------------------------------------------------------------------

\subsection{Universo de activos}

El universo de inversión está compuesto por 30 acciones del S\&P 500,
seleccionadas de modo que cubren los once sectores GICS. La
Tabla~\ref{tab:universo} detalla los tickers y sus sectores.

\begin{table}[htbp]
    \centering
    \caption{Universo de 30 activos por sector GICS.}
    \label{tab:universo}
    \begin{tabular}{ll}
        \toprule
        \textbf{Sector} & \textbf{Tickers} \\
        \midrule
        Tecnología & AAPL, MSFT, ORCL, CSCO, INTC \\
        Salud & JNJ, PFE, MRK, UNH \\
        Finanzas & JPM, GS, AXP, BAC \\
        Consumo discrecional & HD, MCD, NKE \\
        Consumo básico & PG, KO, WMT \\
        Comunicaciones & VZ, T, DIS \\
        Energía & XOM, CVX \\
        Industriales & CAT, HON, UPS \\
        Materiales & FCX \\
        Utilities & NEE \\
        Inmobiliario & AMT \\
        \bottomrule
    \end{tabular}
\end{table}

Como referencia de mercado se utiliza el \textbf{SPY} (SPDR S\&P 500 ETF
Trust), que replica el S\&P 500 y no forma parte del universo a optimizar.

\subsection{Periodo y frecuencia}

\begin{itemize}
    \item \textbf{Periodo total:} 1 de enero de 2010 -- 31 de diciembre de
          2024 (15 años).
    \item \textbf{Datos diarios:} 3.773 días de negociación para 31 tickers
          (30 acciones + SPY).
    \item \textbf{Frecuencia de trabajo:} mensual. Los retornos logarítmicos
          diarios se agregan a frecuencia mensual mediante suma:
          $r_t^{\text{mensual}} = \sum_{d \in \text{mes}(t)} r_d^{\text{diaria}}$,
          propiedad que deriva de la aditividad de los log-retornos
          (ecuación~\ref{eq:aditividad-log}).
    \item \textbf{Meses resultantes:} 180 meses (enero 2010 -- diciembre 2024).
\end{itemize}

\subsection{Fuente y preprocesado}

Los precios ajustados (cierre) se descargan de \textbf{Yahoo Finance} a
través de la librería \texttt{yfinance} de Python. El preprocesado consiste en:

\begin{enumerate}
    \item \textbf{Cálculo de retornos logarítmicos diarios:}
          $r_t = \ln(P_t / P_{t-1})$.
    \item \textbf{Limpieza:} eliminación de días sin negociación,
          \textit{forward-fill} de datos faltantes.
    \item \textbf{Agregación mensual:} suma de retornos diarios dentro de
          cada mes natural.
\end{enumerate}

Se reconoce explícitamente la presencia de \textbf{sesgo de supervivencia}:
el universo está compuesto por constituyentes actuales (2024) del S\&P 500,
no por los constituyentes históricos. Esto implica que las empresas que
quebraron o fueron excluidas del índice durante el periodo de estudio no
están representadas. El sesgo tiende a sobreestimar ligeramente las
rentabilidades históricas, pero su impacto en un estudio de asignación de
activos como este es limitado, ya que el objetivo es comparar metodologías
de optimización, no estimar rentabilidades absolutas.

%--------------------------------------------------------------------
\section{Esquema de backtesting walk-forward}
%--------------------------------------------------------------------

El experimento sigue un protocolo de backtesting walk-forward con
rebalanceo mensual, que constituye el estándar metodológico en la literatura
de optimización de carteras \cite{DeMiguel2009}.

\subsection{Parámetros del backtesting}

\begin{itemize}
    \item \textbf{Ventana de estimación:} 36 meses (3 años).
    \item \textbf{Frecuencia de rebalanceo:} mensual.
    \item \textbf{Periodo out-of-sample:} enero 2013 -- diciembre 2024
          (144 meses, de los cuales se evalúan 143 por el desfase de la
          predicción).
    \item \textbf{Restricciones:} long-only ($w_i \ge 0$), peso máximo por
          activo $w_i \le 0{,}20$, presupuesto $\sum_i w_i = 1$.
\end{itemize}

\subsection{Procedimiento mensual}

Para cada mes $t$ desde enero de 2013 ($t = 37$) hasta noviembre de 2024
($t = 179$): se toma la ventana de 36 meses ($t-35, \dots, t$); se estima
$\boldsymbol{\Sigma}$ con uno de los tres estimadores (muestral, Ledoit-Wolf o
PCA al 80\% de varianza) y $\boldsymbol{\mu}$ con uno de los cuatro métodos
(media histórica, Ridge, Lasso o RF); se resuelve el QP de mínimo riesgo
(ecuación~\ref{eq:qp-riesgo}) barriendo el objetivo de rentabilidad $\mu^*$ y
se toma el punto de
máximo Sharpe, que da los pesos $\mathbf{w}_t$; y se evalúa la rentabilidad
realizada $r_{p,t+1} = \mathbf{w}_t^T \mathbf{r}_{t+1}$.

Esquemáticamente, para un total de $T=180$ meses:
\[
\underbrace{r_2, \dots, r_{37}}_{\text{Train}_1} \to
\underbrace{r_{38}}_{\text{Pred}_1} \quad
\underbrace{r_3, \dots, r_{38}}_{\text{Train}_2} \to
\underbrace{r_{39}}_{\text{Pred}_2} \quad
\cdots \quad
\underbrace{r_{144}, \dots, r_{179}}_{\text{Train}_{143}} \to
\underbrace{r_{180}}_{\text{Pred}_{143}}
\]

donde Train$_k$ es la ventana de 36 meses usada para estimar
$\boldsymbol{\mu}$ y $\boldsymbol{\Sigma}$ en la iteración $k$, y Pred$_k$
es el mes out-of-sample cuya rentabilidad $r_{p,k}$ se evalúa.

\subsection{Prevención del look-ahead}

Para evitar el look-ahead se toman las siguientes medidas:
\begin{itemize}
    \item La estimación de $\boldsymbol{\Sigma}$ en el mes $t$ solo utiliza
          rentabilidades hasta $t$.
    \item Las features para la predicción ML en $t$ solo utilizan datos
          hasta $t$ (ver Capítulo~\ref{ch:ml}).
    \item La rentabilidad out-of-sample $r_{p,t+1}$ \textbf{no} se utiliza
          en ningún paso de la optimización para el mes $t$.
    \item La validación cruzada para la selección de hiperparámetros de los
          modelos ML se realiza con \texttt{TimeSeriesSplit}, que respeta el
          orden cronológico.
\end{itemize}

%--------------------------------------------------------------------
\section{Carteras evaluadas}
%--------------------------------------------------------------------

Se evalúan 8 carteras, presentadas en la Tabla~\ref{tab:carteras-evaluadas}.
Las tres primeras corresponden al modelo de Markowitz clásico con distintos
estimadores de $\boldsymbol{\Sigma}$, las tres siguientes incorporan
predicciones ML de $\boldsymbol{\mu}$, y las dos últimas son baselines
pasivos.

\begin{table}[htbp]
    \centering
    \caption{Carteras evaluadas en el experimento.}
    \label{tab:carteras-evaluadas}
    \begin{tabular}{lll}
        \toprule
        \textbf{Cartera} & \textbf{Estimador de $\boldsymbol{\Sigma}$} & \textbf{Estimador de $\boldsymbol{\mu}$} \\
        \midrule
        Markowitz (muestral) & Muestral & Media histórica \\
        Markowitz (Ledoit-Wolf) & Ledoit-Wolf & Media histórica \\
        Markowitz (PCA) & PCA (80\% var.) & Media histórica \\
        Markowitz + Ridge & Ledoit-Wolf & Ridge \\
        Markowitz + Lasso & Ledoit-Wolf & Lasso \\
        Markowitz + RF & Ledoit-Wolf & Random Forest \\
        1/N (equiponderada) & -- & -- \\
        SPY (mercado) & -- & -- \\
        \bottomrule
    \end{tabular}
\end{table}

\textbf{Nota sobre la elección de $\boldsymbol{\Sigma}$ para las estrategias
ML:} se utiliza el estimador de Ledoit-Wolf para todas ellas, por ser el más
fiable en términos de diversificación y estabilidad
(Capítulo~\ref{ch:limitaciones}). Esto permite aislar el efecto de la mejora
en $\boldsymbol{\mu}$ sin que las diferencias en $\boldsymbol{\Sigma}$
contaminen la comparación.

%--------------------------------------------------------------------
\section{Métricas de evaluación}
%--------------------------------------------------------------------

\subsection{Métricas de rentabilidad-riesgo}

Sobre la serie de 143 rentabilidades mensuales out-of-sample de cada
cartera, se calculan:

\begin{itemize}
    \item \textbf{Rentabilidad anualizada} (compuesta geométricamente):
          \begin{equation}
              r_{\text{anual}} = \Bigl(\prod_{t=1}^{T} (1 + r_{p,t})
              \Bigr)^{12/T} - 1.
              \label{eq:rent-anual}
          \end{equation}
    \item \textbf{Volatilidad anualizada:}
          $\sigma_{\text{anual}} = \sigma_{\text{mensual}} \cdot \sqrt{12}$.
    \item \textbf{Ratio de Sharpe} (anualizado, $r_f = 0$):
          $S_p = r_{\text{anual}} / \sigma_{\text{anual}}$.
    \item \textbf{Máximo drawdown:} máxima caída desde un pico anterior en la
          curva de valor acumulado. Expresado como:
          \begin{equation}
              \text{MDD} = \max_{t} \left(
              1 - \frac{V_t}{\max_{s \le t} V_s}
              \right),
              \label{eq:mdd}
          \end{equation}
          donde $V_t = \prod_{s=1}^{t} (1 + r_{p,s})$ es el valor acumulado
          de la cartera en el mes $t$.
\end{itemize}

\subsection{Métricas de composición de cartera}

\begin{itemize}
    \item \textbf{Turnover medio:} mide la rotación de la cartera entre
          rebalanceos consecutivos,
          \begin{equation}
              \text{TO} = \frac{1}{2(T-1)} \sum_{t=2}^{T}
              \sum_{i=1}^{n} |w_{i,t} - w_{i,t-1}|.
              \label{eq:turnover}
          \end{equation}
          El factor $1/2$ normaliza el turnover entre 0 (cartera estática)
          y 1 (rotación completa). Un turnover elevado implica mayores costes
          de transacción.
    \item \textbf{Índice Herfindahl-Hirschman (HHI) medio:}
          $\text{HHI} = \frac{1}{T} \sum_t \sum_i w_{i,t}^2$.
          Mide la concentración de la cartera: $1/n$ para la equiponderada,
          $1$ para un solo activo.
    \item \textbf{Número medio de activos en cartera:} media del número de
          activos con peso $> 0{,}1\%$ en cada rebalanceo.
\end{itemize}

\subsection{Métricas de predicción ML}

Para los modelos de ML se reportan además el RMSE, el MAE y el $R^2$
out-of-sample de las predicciones de $\boldsymbol{\mu}$ (definidos en el
Capítulo~\ref{ch:ml}); un $R^2$ negativo indica que el modelo predice peor que
la media histórica, aunque la cartera podría aun así mejorar a la clásica si
acierta el orden (ranking) de los activos.

%--------------------------------------------------------------------
\section{Implementación y reproducibilidad}
%--------------------------------------------------------------------

El código está implementado en Python 3.13 y organizado en módulos
independientes bajo \texttt{src/}:

\begin{itemize}
    \item \texttt{data\_ingest/}: descarga y preprocesado de datos.
    \item \texttt{optimization/}: optimización QP (\texttt{cvxpy}),
          estimación de covarianza (\texttt{scikit-learn}) y restricciones.
    \item \texttt{ml/}: feature engineering, modelos (Ridge, Lasso, RF) y
          pipeline de predicción.
    \item \texttt{backtest/}: motor de backtesting walk-forward y métricas
          de evaluación.
\end{itemize}

La carpeta \texttt{tests/} contiene los tests que validan la implementación:
la coincidencia de la solución numérica del QP con la solución analítica y el
cumplimiento de las condiciones KKT (Capítulo~\ref{ch:formulacion}). Se incluye
además, en \texttt{demo/}, una demostración web interactiva del optimizador,
como material de apoyo a la presentación del trabajo.

La reproducibilidad está garantizada mediante:
\begin{itemize}
    \item Semilla fija (\texttt{seed = 42}) para todos los componentes
          aleatorios.
    \item Versiones de dependencias congeladas en \texttt{requirements.txt}.
    \item Datos descargables automáticamente desde Yahoo Finance.
    \item Scripts de ejecución en \texttt{scripts/} que generan todas las
          tablas y figuras de la memoria.
\end{itemize}

%--------------------------------------------------------------------
\section{Síntesis}

El diseño experimental combina 6 carteras optimizadas (tres clásicas, una por
estimador de $\boldsymbol{\Sigma}$, y tres con $\boldsymbol{\mu}$ predicho por
ML sobre Ledoit-Wolf) y 2 baselines pasivos, un backtesting walk-forward de 143 meses out-of-sample (ventanas de 36
meses, rebalanceo mensual) y métricas de rentabilidad-riesgo, composición y
calidad predictiva. El Capítulo~\ref{ch:resultados} presenta los resultados y
el Capítulo~\ref{ch:conclusiones}, las conclusiones.
